\section{Introduction}\label{sec:intro}

A quantum channel is usually described by what it does to one input.
Suppose instead that a device must apply the same channel $n$ times, and
that each output must leave before the next input arrives. With a fresh
pure environment for every use, the device needs almost no memory. But
fresh pure qubits are a resource~\cite{RybarZiman2008}, and a device
given little purity must keep what it cannot afford to throw away. How
much must it keep, and what decides the answer?

We count three resources. The memory $Q$ is the number of qubits the
device keeps between uses. The exchange $J$ is the number of qubits it
hands over for good and replaces with fresh ones. The purity of a
register is its number of qubits minus its
entropy~\cite{HorodeckiHorodeckiOppenheim2003}; $P$ is the purity of
the initial memory and $C$ the total purity of the fresh qubits. Every
reset, purifier and random seed the device keeps is counted. Unless we
say otherwise, fresh qubits are maximally mixed, so $C=0$. The observer
may keep a reference and may choose each input after seeing earlier
outputs. The error is the largest half trace distance, over all such
observers, between this whole experiment and $n$ independent uses of the
channel. A device that releases each output before the next input
arrives is a \emph{causal service}.

The paper has three messages. Immediate release has a price, which a
single qutrit channel already shows. One closed set of channels, those
that a finite maximally mixed bath can imitate, separates cheap from
expensive repeated use. Purity is known to mark this set, and we show
that memory marks it too, for channels with a flat probe. And for
channels built from group relations, at high precision, the memory of
$n$ uses is governed by a single use.

\paragraph{The price of immediacy.}
Consider the Werner--Holevo channel~\cite{WernerHolevo2002} on
$3\times3$ matrices,
\[
 \Theta(X)=\tfrac12(\Tr(X)I-X^T),\qquad H=\log_2 3.
\]
Up to a change of basis it is the spin-one channel that Landau and
Streater used to show that a unital channel need not be a mixture of
unitaries~\cite{LandauStreater1993}.

\begin{introthm}[Theorems~\ref{thm:qutrit-purity} and~\ref{thm:qutrit-collective}]\label{thm:intro-qutrit}
Every causal service of $\Theta$ whose error tends to zero needs
supplied purity $P+C\geq(H-o(1))n$~\cite[Theorem~3.2(c) and Proposition~F.1(c)]{DouglasMghirbiB}, and exact causal
services attain $P+C=nH+O(\sqrt n)$. If all $n$ outputs may be released
together, $\Theta^{\otimes n}$ has an implementation with error
$1/\mathrm{poly}(n)$ and supplied purity $O(\log n)$, refining the rate
zero of~\cite{FaistBertaBrandao2019}.
\end{introthm}

The rate $H$ is the purity cost of an extreme channel, from earlier
work~\cite{DouglasMghirbiB}. What is added here is a bound at finite
error and the comparison with collective delivery. The reason for the
rate is rigidity. The Choi state of $\Theta$ lives on a
three-dimensional antisymmetric subspace, where tracing out the output
acts as $\Theta$ itself, which has a bounded inverse. So a nearly correct
use cannot correlate its output and reference with a purification of
the memory, and each use must deposit $\log_2 3$ fresh bits of entropy
there. Collective delivery avoids this by holding all inputs until the
end, so its memory is linear.

Letting outputs wait interpolates between the two. At error $n^{-2}$,
with sublinear memory and linear exchange, batches much shorter than
$\log n$ keep the full rate, while batches much longer than $\log n$
have rate zero (Theorem~\ref{thm:batch-purity}). Rate zero for batches
much longer than $\log n$ holds for every unital channel
(Theorem~\ref{thm:unital-batch}). The same logarithmic threshold appears
for a closed apparatus delivering an extreme ququart Schur
channel~\cite[Theorems~7 and~8]{DouglasSchur}; we carry it to open
services.

The collective bound holds for every unital channel
(Lemma~\ref{lem:collective-gaussian}, which
extends~\cite[Theorem~1]{DouglasSchur}). At the level of rates it is a
theorem of Faist, Berta and Brand\~ao. With trivial Hamiltonians,
thermal operations are the noisy operations, work is counted in pure
qubits, and unital channels have zero thermodynamic
capacity~\cite{FaistBertaBrandao2019},
\cite[Equation~(4.9) and Theorem~7.1]{FaistBertaBrandao2021}. The lemma
adds the explicit $O(\log n)$ at polynomially small error. Classical
randomness cannot explain it. The channel $\Theta$ is not
factorizable~\cite[Example~3.1]{HaagerupMusat2011}, so its tensor powers
stay a fixed distance from mixtures of
unitaries~\cite[Section~6]{HaagerupMusat2011}. This refuted the
asymptotic quantum Birkhoff conjecture of Smolin, Verstraete and
Winter~\cite{SmolinVerstraeteWinter2005}; see also~\cite{Shor2010}.

\paragraph{The dividing line.}
A \emph{finite tracial factorization} of a channel is an implementation
by one unitary acting on the input and a finite bath that starts
maximally mixed. It uses no purity, and on a maximally mixed input it
produces no bath entropy. Horodecki, Horodecki and Oppenheim introduced
these implementations as noisy
operations~\cite{HorodeckiHorodeckiOppenheim2003}, and Gour et al.\
define noisy operations as their closure~\cite{GourEtAl2015}. Let
$\overline{\mathcal F}$ be this closure. By Haagerup and Musat,
$\overline{\mathcal F}$ consists of the factorizable channels that admit
an ancilla embeddable in an ultrapower of the hyperfinite
$\mathrm{II}_1$ factor~\cite[Theorem~3.6]{HaagerupMusat2015}, and it
contains every factorizable channel exactly when Connes' embedding
problem has a positive answer~\cite[Theorem~3.7]{HaagerupMusat2015}.
Distance is the normalized Choi distance $d_J$, half the trace distance
between Choi states.

A channel of Choi rank $r$ has a \emph{flat probe} if some input, with a
reference, leaves its minimal environment maximally mixed; equivalently,
its entropy exchange~\cite{Schumacher1996} can reach $\log_2 r$. For such
a channel in $\overline{\mathcal F}$, about $\frac12\log r$ exchanged
qubits per use is the least that permits sublinear memory as the error
tends to zero~\cite[Theorem~3.4]{DouglasMghirbiB}. From here on, unless
we say otherwise, the exchange is at most $\lfloor(n-1)\log r/2\rfloor$
and errors are at most $\frac1{16}$.

\begin{introthm}[Theorem~\ref{thm:closure-law}]\label{thm:intro-closure}
Let $\Phi$ be a channel on $M_d$.
\begin{enumerate}[label=(\alph*),itemsep=1pt]
\item If $\Phi\in\overline{\mathcal F}$ has a flat probe of rank
$r\geq2$, then for every error sequence $\epsilon_n>0$ with
$\log(1/\epsilon_n)=o(n)$ there are causal services with error at most
$\epsilon_n$ and memory $Q=o(n)$.
\item \cite[Theorem~3.2(b)]{DouglasMghirbiB} If $\Phi$ is at normalized
Choi distance $\Delta>0$ from $\overline{\mathcal F}$, every causal
service with imports between uses and error $\epsilon$ has, at any
exchange, $P+C\geq{n(\Delta-\epsilon)_+^2/(2\ln2)}$. With maximally
mixed imports its memory is therefore linear whenever $\epsilon<\Delta$.
\end{enumerate}
\end{introthm}

Part~(a) sharpens~\cite[Corollary~3.3(c)]{DouglasMghirbiB}, which gives
sublinear memory at an error that tends to zero at no prescribed rate.
The mechanism of part~(b) is local
tracialization~\cite[Lemma~5.4]{DouglasMghirbiB}: a use that raises the
bath entropy only slightly, on a maximally mixed input, is close in
Choi distance to a mixture of finite tracial factorizations
(Lemma~\ref{lem:tracialization}). A channel a fixed distance from
$\overline{\mathcal F}$ must therefore produce bath entropy at a
positive rate, and supplied purity pays for it. The qutrit channel
$\Theta$ lies at a positive distance from $\overline{\mathcal F}$,
computed by Haagerup and Musat~\cite[Theorem~5.6 and
Corollary~5.7]{HaagerupMusat2015} and restated in Choi distance in
Proposition~\ref{prop:qutrit-gap}.

The question of reusing one device without resets goes back to Ryb\'ar
and Ziman~\cite{RybarZiman2008}. In their model one fixed unitary acts
on a finite memory and each input, and a channel is repeatable if every
use, on uncorrelated inputs, induces exactly that channel. They showed
that random-unitary channels are repeatable and that nonunital channels
are not repeatable with a finite memory. They left the other unital
channels open~\cite[Section~IV]{RybarZiman2008}. Their entropy
count~\cite[Equation~(3.7)]{RybarZiman2008} is the prototype of the
purity bookkeeping behind Theorem~\ref{thm:intro-closure}(b).

Proposition~\ref{prop:repeatable} answers their question: some unital
channels that are not random-unitary are repeatable, and some unital
channels are not. The positive half rests on known results. A Schur
channel of Haagerup and Musat~\cite[Example~3.3]{HaagerupMusat2011} is
unital and not random-unitary. Lie, Son, Boes, Ng and Wilming showed
that it has a catalytic dilation, in which the memory state never
changes~\cite[Proposition~15 and Appendix~C]{LieSonBoesNgWilming2026}.
Such a device repeats its channel, here with one memory qubit. For the
negative half we use the same bookkeeping. A device that repeats $\Phi$
for $n$ uses with memory dimension $M$ has
$\log M\geq n\operatorname{dist}_J(\Phi,\overline{\mathcal F})^2/(2\ln2)$.
A channel that one finite memory repeats forever is a finite mixture of
finite tracial factorizations. The qutrit channel $\Theta$ is unital, but
repeating it $n$ times needs $\log M\geq n\log3$, which their $n$-cell
device attains.

Ryb\'ar and Ziman also proposed devices whose memory starts maximally
mixed, which are our finite tracial factorizations, followed by a reset
to the maximally mixed state. A reused reset device receives inputs
correlated with its own memory, and repeatability says nothing about
such inputs. Theorem~\ref{thm:intro-closure}(a) instead replaces about
$\frac12\log r$ memory qubits per use by fresh maximally mixed ones,
under joint error and with memory that may grow sublinearly. It can be
read as an open-system, quantitative form of their proposal. Lie, Son,
Boes, Ng and Wilming call the finite mixtures of finite tracial
factorizations strongly factorizable~\cite{LieSonBoesNgWilming2026};
their closure is $\overline{\mathcal F}$.

Exact factorizability and membership in $\overline{\mathcal F}$ are
different questions. Haagerup and Musat's equivalence, combined with
$\mathrm{MIP}^*=\mathrm{RE}$~\cite{JiEtAl2020}, implies nonconstructively
that in some finite dimension an exactly factorizable channel lies
outside $\overline{\mathcal F}$; Musat records this for a Schur
channel~\cite[Corollary~3.11]{Musat2025}. Our companion
paper~\cite{MghirbiDouglasD} gives explicit rational examples, conditional
on recent preprints. Their distance from $\overline{\mathcal F}$ is explicit
but astronomically small (Proposition~\ref{prop:rational-distances}). In the other direction, some
mixtures of two unitaries lie in $\overline{\mathcal F}$ but have no
finite tracial factorization~\cite[Proposition~3.6]{Musat2025}. Musat
and R\o rdam give channels in $\overline{\mathcal F}$ with no
finite-dimensional ancilla at all, in every dimension
$d\geq11$~\cite[Theorem~4.1]{MusatRordam2019}; Mori lowers this to
$d\geq8$~\cite[Theorem~2.5]{Musat2025}.

\paragraph{Three regimes for group gadgets.}
A \emph{canonical gadget} (Section~\ref{sec:gadgets}) is a channel built
from a finite list of relations in a group, by the construction
of~\cite[Section~9]{DouglasMghirbiC}. Its Kraus operators are indexed by
finitely many group elements, the \emph{labels}, and it is factorizable
through the group von Neumann algebra. Write $P_\Phi$ for the group
presented by the gadget's own relations. For gadgets the dividing line
refines into three regimes at every fixed error
(Theorem~\ref{thm:trichotomy} and Table~\ref{tab:regimes}). Memory is at
most logarithmic when a finite quotient of $P_\Phi$ separates the
labels; the gadget then has an exact finite tracial
factorization~\cite[Theorem~8.6]{DouglasMghirbiB}. Memory diverges but
stays sublinear when the labels are separated only by unitaries in a
tracial ultrapower of the hyperfinite factor. Otherwise it is linear.
The linear regime contains a gadget exactly when some finitely presented
group is not hyperlinear (Proposition~\ref{prop:regime-iii}).

Slofstra and Vidick read the same three-way split off hyperlinear
profiles~\cite[Section~2]{SlofstraVidick2018}. For linear-system games
it governs the entanglement needed for near-perfect
play~\cite{CleveLiuSlofstra2017,Slofstra2019,SlofstraVidick2018}. There
the dimension diverges as the error tends to zero; here memory diverges
with the number of uses at fixed error.

\paragraph{One use governs many.}
Realize one use of a gadget $\Phi$ by a unitary acting on the input and a
bath in state $\rho$. Over $n$ uses the bath entropy $S=S(\rho)$ is paid
once, and the entropy $a$ that one use adds to the bath (with maximally
mixed input) is paid $n$ times. So the natural cost of the use is
$S+na$. Let the \emph{profile} $Z_\Phi(n)$ be the least such cost over
uses whose output stays exactly in the span of the gadget's Kraus
operators and which tell the labels apart to a fixed coarse accuracy
(Definition~\ref{def:exact-profile}). The bath state may be coherent.

\begin{introthm}[Theorem~\ref{thm:main}]\label{thm:intro-memory}
Let $\Phi$ be a canonical gadget with Choi rank $r\geq2$, and let
$\epsilon\leq1/[n\log_2^2(en)]$. The least memory $Q$ of a causal service
of $\Phi$ with error $\epsilon$ satisfies
\[
 c\,(1+Z_\Phi(n))\ \leq\ 1+Q\ \leq\ C\log_2^2(n/\epsilon)\,(1+Z_\Phi(n)),
\]
with constants that depend only on the dimension and the Choi rank.
\end{introthm}

At fixed error the lower bound is a second profile, $L^*_\Phi(n)$, which
also allows each use a leakage of order $1/n$ outside the Kraus span
(Definition~\ref{def:profile}). Whether $Z_\Phi(n)$ and $L^*_\Phi(n)$
agree up to polylogarithmic factors is the main open problem
(Problem~\ref{prob:keep-bias}).

The profiles make intermediate scales visible. For the configuration-lamp
gadget of~\cite[Section~11]{DouglasMghirbiC}, memory is
$n^{1/2+o(1)}$~\cite[Theorem~11.1(b)]{DouglasMghirbiC}, and the
fixed-error profile $L^*_\Phi$ matches it (Corollary~\ref{cor:lamp}).
For Houghton's channel, devices whose uses act through Houghton's group,
with negligible discard, need memory $n^{1/3}$ up to logarithmic
factors, and such devices attain it
(Theorem~\ref{thm:houghton-genuine}). For arbitrary devices, $n^{1/3}$
up to logarithms is an upper bound~\cite[Corollary~6.7]{DouglasMghirbiB},
and the exponent is open.

\paragraph{Related work.}
Memory channels~\cite{KretschmannWerner2005,CarusoEtAl2014} and collision
models~\cite{CiccarelloEtAl2022} describe repeated interactions with an
environment. Sequential protocols are described by quantum
combs~\cite{ChiribellaDArianoPerinotti2009} and
strategies~\cite{GutoskiWatrous2007}; our error uses their adaptive
discrimination setting. The memory needed to realize a given comb has
been minimized when discarding is free~\cite{BisioEtAl2011,BisioEtAl2012};
here discarding costs exchange and every retained qubit counts.
Programmable processors store a channel in a program
register~\cite{NielsenChuang1997,HilleryZimanBuzek2002}; here the channel
is fixed and the memory serves its repetition. Process tensors describe
the multi-time response to interventions and detect environmental
memory~\cite{PollockEtAl2018}. Here the target process has independent
uses, and the memory belongs to its resource-constrained implementation.
Quantum statistical complexity minimizes simulator entropy for
prescribed output statistics~\cite{GuEtAl2012}, not for a channel on
arbitrary entangled inputs. Mendl and Wolf analyze the mixed-unitary
structure of unital and covariant channels~\cite{MendlWolf2009}; Haagerup
and Musat determined which members of their qutrit family are
factorizable~\cite[Corollary~5.7]{HaagerupMusat2015}, the input to
Proposition~\ref{prop:qutrit-gap}. On the operator-algebra side, a
unital channel is factorizable exactly when its iterates are
compressions of the iterates of one trace-preserving automorphism of a
larger algebra, equivalently when it has a Markov
dilation~\cite[Theorem~4.4]{HaagerupMusat2011}; Haagerup and Musat
credit the first equivalence to K\"ummerer. These are iterates of the
channel, not independent uses.

Thermodynamics prices processes in work. Faist and Renner give the
single-shot work cost of a process on a specified
input~\cite{FaistRenner2018}. Faist, Dupuis, Oppenheim and Renner bound
the work cost of a logical process on a fixed input by a smooth
max-entropy of the discarded information given the
output~\cite{FaistDupuisOppenheimRenner2015}. Both allow every
Gibbs-preserving map for free; with trivial Hamiltonians these are the
unital channels, so $\Theta$ costs nothing there. Gibbs-preserving maps
can outperform thermal operations on single
instances~\cite{FaistOppenheimRenner2015}, but for i.i.d.\
implementations the two cost the same~\cite{FaistBertaBrandao2019},
\cite[Theorem~7.1]{FaistBertaBrandao2021}.
Theorem~\ref{thm:intro-qutrit} shows that the gap reopens, at the full
rate $H$, when each output must leave at once. Faist, Berta and
Brand\~ao asked how their costs change under distinguishability measures
suited to memory effects~\cite{FaistBertaBrandao2019}. For trivial
Hamiltonians the answer depends on the delivery contract: the qutrit
lower bound needs only fresh entangled inputs chosen in advance, so the
jump comes from immediate release, not from the stronger measure.

Catalysis is the closest neighbour. Exact return of a catalyst's marginal
can coexist with correlations: entropy and rank govern unitary state
conversion~\cite{Wilming2022}, and free energy governs approximate
correlated-catalytic conversion under Gibbs-preserving
operations~\cite{ShiraishiSagawa2021}. Lie and Jeong show that every
catalytic channel is factorizable, so that not every unital channel is
catalytic, and treat uncorrelatedness as a consumable randomness
resource~\cite{LieJeong2021}. Boes et al.\ show that a quantum source of
randomness of dimension $\sqrt d$ suffices for every noisy state
transition on a $d$-dimensional system~\cite{BoesEtAl2018}, a halving of
the same kind as the exchange rate $\frac12\log r$. Son et al.\ show that
exact catalyst return near one input implies return on every
input~\cite{SonEtAl2026}. Returning a marginal is weaker than simulating
a full adaptive experiment; here retained catalysts count as memory and
every import is independent and charged. Channel resource theories
quantify processes under specified free
operations~\cite{GourWinter2019}, thermodynamic ones
included~\cite{FaistBertaBrandao2021}; our additional constraint is the
chronological memory and exchange budget.

\paragraph{Relation to earlier papers.}
The Houghton paper~\cite{DouglasMghirbiA} supplies the far-commutator
area bound and the finite models of Houghton's group; the presentation
itself is Johnson's, as used in~\cite[Section~4]{DouglasMghirbiB}. The
causal-memory paper~\cite{DouglasMghirbiB} supplies the model, the step
from a service to a single good use, the service at the optimal exchange
rate, and the purity bounds behind Theorem~\ref{thm:intro-qutrit} and
Theorem~\ref{thm:intro-closure}(b). The configuration-lamps
paper~\cite{DouglasMghirbiC} supplies the gadget construction and the
lamp channel's memory laws. The rational-channel
paper~\cite{MghirbiDouglasD} supplies exactly factorizable channels at an
explicit distance from $\overline{\mathcal F}$, and the group certificate
behind our canonical example in regime~(iii), conditional on recent
preprints~\cite{WangZhi2026,AlekseevLiuThom2026,Thom2026}. Closed repeated
use~\cite{DouglasClosedMemory}, the no-reset setting
of~\cite{RybarZiman2008}, concerns one unchanging bath with no exchange.
Finite numerical checks of the constants are in the accompanying
\texttt{verify} directory; they support but do not replace the proofs.

\paragraph{Organization.}
Section~\ref{sec:model} specifies the model. The qutrit separation comes
first, followed by the dividing line and the one-use memory bounds. The
remaining sections develop examples and the open comparison. Each main
statement links to its proof.

The appendices fall into three groups. Appendices~\ref{app:purity},
\ref{app:closure-proofs}, \ref{app:compactness}, \ref{app:graded},
\ref{app:static} and~\ref{app:lower} prove the results of
Sections~\ref{sec:model}--\ref{sec:lower}: the qutrit and batching
results, the closure law and repeatability, the gadget classification,
the memory law with its static comparisons, and the one-round lower
bound. Appendix~\ref{app:houghton} proves the Houghton results of
Section~\ref{sec:houghton}, which also use the device of
Theorem~\ref{thm:houghton-groups} and the cap of
Proposition~\ref{prop:sixth}. The rest is specialized.
Appendix~\ref{sec:monomial} bounds a mixture profile for phased monomial
dynamics. Appendices~\ref{sec:flat}, \ref{sec:biased}
and~\ref{sec:obstructions}, with proofs in Appendices~\ref{app:frames},
\ref{app:flat}, \ref{app:biased} and~\ref{app:obstructions}, develop
flat-source, conserved-bias and spatial constructions and their limits.
Appendices~\ref{app:houghton-sources}--\ref{app:houghton-moderate} extend
the Houghton analysis; Section~\ref{sec:houghton} summarizes their
results.
